\documentclass[12pt,a4paper]{article}

% -------------------------------------------------
% PACOTES E CONFIGURAÇÕES
% -------------------------------------------------
\usepackage[utf8]{inputenc}
\usepackage[T1]{fontenc}

% Suporte a idioma com fallback automático para compilação em qualquer ambiente TeX Live
\IfFileExists{brazilian.ldf}{%
  \usepackage[brazilian]{babel}%
}{%
  \usepackage[english]{babel}%
  \addto\captionsenglish{%
    \renewcommand{\contentsname}{Sumário}%
    \renewcommand{\listfigurename}{Lista de Figuras}%
    \renewcommand{\listtablename}{Lista de Tabelas}%
    \renewcommand{\figurename}{Figura}%
    \renewcommand{\tablename}{Tabela}%
  }%
}

\usepackage{lmodern}
\usepackage{geometry}
\geometry{top=2.8cm,bottom=2.5cm,left=3cm,right=2.5cm}

% Microtype: melhora tipográfica (carregado se disponível)
\IfFileExists{microtype.sty}{\usepackage{microtype}}{}

\usepackage{setspace}
\onehalfspacing
\usepackage{graphicx}
\usepackage{float}
\usepackage{booktabs}
\usepackage{longtable}
\usepackage{array}
\usepackage{tabularx}

% Enumitem com fallback elegante
\IfFileExists{enumitem.sty}{%
  \usepackage{enumitem}%
  \setlist[enumerate,1]{label=\textbf{\arabic*.}, leftmargin=*}%
}{%
  \renewcommand{\labelenumi}{\textbf{\arabic{enumi}.}}%
}

\usepackage{xcolor}
\usepackage{hyperref}
\usepackage{bookmark}
\usepackage{listings}
\usepackage{amsmath}
\usepackage{amssymb}
\usepackage{fancyhdr}

% Lastpage com fallback seguro
\IfFileExists{lastpage.sty}{%
  \usepackage{lastpage}%
}{%
  \AtEndDocument{\label{LastPage}}%
}

\usepackage{etoolbox}
\renewcommand{\_}{\textunderscore\allowbreak}

% -------------------------------------------------
% CORES DO PROJETO LCQUI
% -------------------------------------------------
\definecolor{entityHead}{RGB}{30,64,110}
\definecolor{entityBg}{RGB}{235,242,250}
\definecolor{entityFrame}{RGB}{30,64,110}

\definecolor{noteHead}{RGB}{176,98,10}
\definecolor{noteBg}{RGB}{255,244,229}
\definecolor{noteFrame}{RGB}{176,98,10}

\definecolor{ruleHead}{RGB}{20,110,80}
\definecolor{ruleBg}{RGB}{230,247,240}
\definecolor{ruleFrame}{RGB}{20,110,80}

\definecolor{pkColor}{RGB}{176,40,40}
\definecolor{fkColor}{RGB}{30,90,160}
\definecolor{typeColor}{RGB}{60,60,60}
\definecolor{typeBg}{RGB}{225,225,232}

% -------------------------------------------------
% HYPERSETUP
% -------------------------------------------------
\hypersetup{
    colorlinks=true,
    linkcolor=blue!50!black,
    urlcolor=blue!60!black,
    citecolor=green!40!black,
    pdftitle={Projeto de Desenvolvimento do Sistema para o Nucleo de Situação de Saúde},
    pdfauthor={Zadoque Carneiro}
}

\pagestyle{fancy}
\fancyhf{}
\lhead{Núcleo de Situação de Saúde}
\rhead{Planejamento Técnico}
\cfoot{\thepage\ de \pageref{LastPage}}
\setlength{\headheight}{15pt}

% -------------------------------------------------
% DEFINIÇÃO DE LINGUAGENS (MERMAID E TYPESCRIPT)
% -------------------------------------------------
\lstdefinelanguage{Mermaid}{%
  morekeywords={flowchart,graph,sequenceDiagram,classDiagram,erDiagram,stateDiagram,participant,subgraph,end,style,click,Note,rect,alt,else,loop,par,and,autonumber},%
  sensitive=false%
}

\lstdefinelanguage{TypeScript}{
  keywords={typeof, new, true, false, catch, function, return, null, switch, var, if, in, while, do, else, case, break, const, let, export, import, from, interface, type, class, extends, implements, async, await, string, number, boolean, any, void, as},
  keywordstyle=\color{blue}\bfseries,
  ndkeywords={admin, onCall, PDFDocument, Date, Promise, reduce},
  ndkeywordstyle=\color{purple}\bfseries,
  identifierstyle=\color{black},
  sensitive=true,
  comment=[l]{//},
  morecomment=[s]{/*}{*/},
  commentstyle=\color{gray}\ttfamily\itshape,
  stringstyle=\color{green!50!black}\ttfamily,
  morestring=[b]',
  morestring=[b]",
  morestring=[b]`
}

\lstset{
    basicstyle=\ttfamily\small,
    breaklines=true,
    frame=single,
    columns=fullflexible,
    keepspaces=true,
    showstringspaces=false,
    literate=
        {á}{{\'a}}1 {à}{{\`a}}1 {ã}{{\~a}}1 {â}{{\^a}}1
        {é}{{\'e}}1 {ê}{{\^e}}1
        {í}{{\'i}}1
        {ó}{{\'o}}1 {õ}{{\~o}}1 {ô}{{\^o}}1
        {ú}{{\'u}}1 {ü}{{\"u}}1
        {ç}{{\c c}}1
        {Á}{{\'A}}1 {À}{{\`A}}1 {Ã}{{\~A}}1 {Â}{{\^A}}1
        {É}{{\'E}}1 {Ê}{{\^E}}1
        {Í}{{\'I}}1
        {Ó}{{\'O}}1 {Õ}{{\~O}}1 {Ô}{{\^O}}1
        {Ú}{{\'U}}1 {Ü}{{\"U}}1
        {Ç}{{\c C}}1
}

\newcommand{\placeholderfigure}[3]{%
\begin{figure}[H]
    \centering
    \fbox{\begin{minipage}[c][#2][c]{0.93\textwidth}
        \centering
        \textbf{Espaço reservado para diagrama/imagem}\\[0.5em]
        \textit{#1.png}\\[0.5em]
        \small #3
    \end{minipage}}
    \caption{#3}
\end{figure}}

% -------------------------------------------------
% CAIXAS DE ENTIDADE, OBSERVAÇÃO E REGRAS
% Suporta tcolorbox quando disponível, ou renderização nativa ultra-rápida
% -------------------------------------------------
\newif\ifusetcolorbox
\IfFileExists{tcolorbox.sty}{\usetcolorboxtrue}{\usetcolorboxfalse}

\ifusetcolorbox
  \usepackage{tcolorbox}
  \tcbuselibrary{skins,breakable}
  \newtcolorbox{entidade}[1]{%
    enhanced, breakable, colback=entityBg, colframe=entityFrame,
    colbacktitle=entityHead, coltitle=white, fonttitle=\bfseries\large,
    title={\ttfamily #1}, boxrule=0.9pt, arc=2mm,
    left=5mm, right=5mm, top=2mm, bottom=2mm, boxsep=1mm,
    attach boxed title to top left={xshift=3mm,yshift*=-1mm},
    boxed title style={boxrule=0pt,arc=1.5mm},
    before skip=10pt, after skip=14pt
  }
  \newcommand{\tipo}[1]{%
    \tcbox[nobeforeafter,boxrule=0pt,boxsep=1pt,left=3pt,right=3pt,top=1.5pt,bottom=1.5pt,
           colback=typeBg,colframe=typeBg,arc=1.5mm,
           fontupper=\scriptsize\ttfamily\color{typeColor}]{#1}%
  }
  \newcommand{\chave}[1]{%
    \ifstrequal{#1}{PK}{\tcbox[nobeforeafter,boxrule=0pt,boxsep=1pt,left=3pt,right=3pt,top=1.5pt,bottom=1.5pt,colback=pkColor,colframe=pkColor,arc=1.5mm,fontupper=\scriptsize\bfseries\color{white}]{PK}}{%
    \ifstrequal{#1}{FK}{\tcbox[nobeforeafter,boxrule=0pt,boxsep=1pt,left=3pt,right=3pt,top=1.5pt,bottom=1.5pt,colback=fkColor,colframe=fkColor,arc=1.5mm,fontupper=\scriptsize\bfseries\color{white}]{FK}}{%
    \tcbox[nobeforeafter,boxrule=0pt,boxsep=1pt,left=3pt,right=3pt,top=1.5pt,bottom=1.5pt,colback=gray!25,colframe=gray!25,arc=1.5mm,fontupper=\scriptsize\itshape\color{black!70}]{#1}%
    }}%
  }
  \newtcolorbox{explicacao}[1][Observação]{%
    enhanced, breakable, colback=noteBg, colframe=noteFrame,
    colbacktitle=noteHead, coltitle=white, fonttitle=\bfseries\small,
    title={\raisebox{-0.1em}{\textbf{i}}\ \ #1}, boxrule=0.8pt, arc=1.5mm,
    left=4mm, right=4mm, top=1.5mm, bottom=1.5mm, before skip=8pt, after skip=10pt
  }
  \newtcolorbox{regra}[1][Regra de negócio]{%
    enhanced, breakable, colback=ruleBg, colframe=ruleFrame,
    colbacktitle=ruleHead, coltitle=white, fonttitle=\bfseries\small,
    title={#1}, boxrule=0.8pt, arc=1.5mm,
    left=4mm, right=4mm, top=1.5mm, bottom=1.5mm, before skip=8pt, after skip=10pt
  }
\else
  % Renderização nativa limpa, rápida e estética (zero dependência de pacotes extras)
  \newcommand{\tipo}[1]{\colorbox{typeBg}{\scriptsize\ttfamily\color{typeColor}\strut\,#1\,}}
  \newcommand{\chave}[1]{%
    \ifstrequal{#1}{PK}{\colorbox{pkColor}{\scriptsize\bfseries\color{white}\strut\,PK\,}}{%
    \ifstrequal{#1}{FK}{\colorbox{fkColor}{\scriptsize\bfseries\color{white}\strut\,FK\,}}{%
    \colorbox{gray!25}{\scriptsize\itshape\color{black!70}\strut\,#1\,}%
    }}%
  }
  \newsavebox{\entidadeBox}
  \newenvironment{entidade}[1]{%
    \par\vspace{8pt}\noindent%
    \def\entidadeTitle{#1}%
    \begin{lrbox}{\entidadeBox}%
    \begin{minipage}{\dimexpr\textwidth-2\fboxrule-10mm\relax}%
    \vspace{2mm}%
  }{%
    \vspace{2mm}%
    \end{minipage}%
    \end{lrbox}%
    \setlength{\fboxrule}{0.9pt}%
    \setlength{\fboxsep}{0pt}%
    \fcolorbox{entityFrame}{entityBg}{%
      \begin{minipage}{\dimexpr\textwidth-2\fboxrule\relax}%
        \colorbox{entityHead}{%
          \begin{minipage}{\dimexpr\textwidth-2\fboxrule-8mm\relax}%
            \vspace{2mm}%
            \hspace{4mm}{\color{white}\bfseries\large\ttfamily \entidadeTitle}%
            \vspace{2mm}%
          \end{minipage}%
        }\par
        \vspace{2mm}%
        \hspace{5mm}\usebox{\entidadeBox}%
        \vspace{2mm}%
      \end{minipage}%
    }\par\vspace{10pt}%
  }
  \newsavebox{\explicacaoBox}
  \newenvironment{explicacao}[1][Observação]{%
    \par\vspace{8pt}\noindent%
    \def\explicacaoTitle{#1}%
    \begin{lrbox}{\explicacaoBox}%
    \begin{minipage}{\dimexpr\textwidth-2\fboxrule-8mm\relax}%
    \small\vspace{2mm}%
  }{%
    \vspace{2mm}%
    \end{minipage}%
    \end{lrbox}%
    \setlength{\fboxrule}{0.8pt}%
    \setlength{\fboxsep}{0pt}%
    \fcolorbox{noteFrame}{noteBg}{%
      \begin{minipage}{\dimexpr\textwidth-2\fboxrule\relax}%
        \colorbox{noteHead}{%
          \begin{minipage}{\dimexpr\textwidth-2\fboxrule-6mm\relax}%
            \vspace{1.5mm}%
            \hspace{3mm}{\color{white}\bfseries\small\raisebox{-0.1em}{\textbf{i}}\ \ \explicacaoTitle}%
            \vspace{1.5mm}%
          \end{minipage}%
        }\par
        \vspace{2mm}%
        \hspace{4mm}\usebox{\explicacaoBox}%
        \vspace{2mm}%
      \end{minipage}%
    }\par\vspace{10pt}%
  }
  \newsavebox{\regraBox}
  \newenvironment{regra}[1][Regra de negócio]{%
    \par\vspace{8pt}\noindent%
    \def\regraTitle{#1}%
    \begin{lrbox}{\regraBox}%
    \begin{minipage}{\dimexpr\textwidth-2\fboxrule-8mm\relax}%
    \small\vspace{2mm}%
  }{%
    \vspace{2mm}%
    \end{minipage}%
    \end{lrbox}%
    \setlength{\fboxrule}{0.8pt}%
    \setlength{\fboxsep}{0pt}%
    \fcolorbox{ruleFrame}{ruleBg}{%
      \begin{minipage}{\dimexpr\textwidth-2\fboxrule\relax}%
        \colorbox{ruleHead}{%
          \begin{minipage}{\dimexpr\textwidth-2\fboxrule-6mm\relax}%
            \vspace{1.5mm}%
            \hspace{3mm}{\color{white}\bfseries\small \regraTitle}%
            \vspace{1.5mm}%
          \end{minipage}%
        }\par
        \vspace{2mm}%
        \hspace{4mm}\usebox{\regraBox}%
        \vspace{2mm}%
      \end{minipage}%
    }\par\vspace{10pt}%
  }
\fi

% Linha de campo: nome | tipo postgres | chave/observação (opcional)
\newcommand{\campo}[3]{%
  \noindent\textit{\textbf{#1}}\hfill\tipo{#2}%
  \ifstrempty{#3}{\ \chave{NOT NULL}}{\ \chave{#3}}%
  \par\vspace{2pt}%
}

\newcommand{\subcampo}[1]{%
  \noindent\hspace{4mm}{\small\textit{#1}}\par\vspace{2pt}%
}

\title{\textbf{Projeto de Desenvolvimento do Sistema para o Núcleo de Situação de Saúde}\\[0.5em]\large Planejamento Funcional, Arquitetural, de Dados e de Implementação}
\author{Zadoque Carneiro}
\date{\today}

\begin{document}
\maketitle
\thispagestyle{empty}
\newpage

\tableofcontents
\newpage

\section{Introdução}
O presente documento detalha o planejamento arquitetural e funcional inicial para o sistema do Núcleo de Situação de Saúde (também referenciado administrativamente como Sala de Situação de Saúde) da Universidade Estadual do Norte Fluminense Darcy Ribeiro (UENF). 

O sistema tem como objetivo principal viabilizar a visualização e análise de dados epidemiológicos provenientes do Sistema de Informação de Agravos de Notificação (SINAN), operando através de uma plataforma web com forte ênfase na interatividade espacial e na clareza das informações.

\section{Arquitetura do Sistema}
Para garantir modularidade, segurança e desempenho, foi adotada uma arquitetura dividida em três camadas principais: Front-end, Back-end Principal (Java) e Serviço Especializado de Dados (Python).

\begin{regra}[Fluxo Central de Comunicação]
O Front-end nunca fará requisições diretas ao serviço Python ou ao Banco de Dados. Todo o fluxo de dados, autenticação e autorização será orquestrado pelo Back-end em Java.
\end{regra}

\subsection{Componentes da Arquitetura}
\begin{enumerate}
    \item \textbf{Front-end (Aplicação Web):} Responsável pela interface do usuário e pelas interações com os mapas.
    \item \textbf{Back-end Principal (Java + Spring Boot):} Atua como o maestro do sistema. Suas responsabilidades incluem:
    \begin{itemize}
        \item Gerenciamento de Autenticação e Autorização de usuários.
        \item Conexão primária e exclusiva com o banco de dados PostgreSQL.
        \item Recepção de requisições do front-end e comunicação restrita com o serviço Python.
    \end{itemize}
    \item \textbf{Serviço de Processamento de Dados (Python):} Microsserviço isolado encarregado especificamente de baixar, processar e tratar os dados do SINAN através da biblioteca \texttt{pySUS}. Este serviço atende única e exclusivamente às chamadas internas feitas pelo servidor Java.
    \item \textbf{Banco de Dados (PostgreSQL):} Responsável pelo armazenamento persistente de credenciais, logs e eventuais consolidações de dados.
\end{enumerate}

\section{Escopo de Interface e Interação Visual}
O escopo inicial do front-end é focado na exploração geográfica dos dados através de mapas interativos, operando no formato de ``drill-down'' (aproximação sucessiva).

\subsection{Regras de Navegação no Mapa}
O mapa interativo possuirá comportamento reativo (hover) em todas as camadas de visualização e seguirá o fluxo abaixo:

\begin{explicacao}[Nível 1: Regiões do Brasil]
O mapa inicial exibe o Brasil subdividido em 5 grandes regiões. Apenas a \textbf{Região Sudeste} possuirá dados reais e estará habilitada para aprofundamento.
\end{explicacao}

\begin{explicacao}[Nível 2: Estados do Sudeste]
Ao clicar na Região Sudeste, a visualização se expande. Os estados da região tornam-se selecionáveis, porém apenas o \textbf{Rio de Janeiro (RJ)} conterá dados alimentados para prosseguir.
\end{explicacao}

\begin{explicacao}[Nível 3: Municípios do Rio de Janeiro]
Ao clicar no RJ, o mapa expande novamente exibindo seus municípios. Nesta versão inicial, a seleção e visualização de dados estarão restritas às seguintes cidades:
\begin{itemize}
    \item Campos dos Goytacazes
    \item Macaé
    \item São João da Barra
    \item Itaperuna
\end{itemize}
\end{explicacao}

\section{Considerações Tecnológicas do Front-end}
O escopo define a tecnologia do Back-end, no entanto, a escolha para o Front-end encontra-se em fase de avaliação estratégica entre desempenho bruto e velocidade de entrega.

\begin{regra}[Decisão Tecnológica do Front-end]
A escolha final deve ponderar a capacidade da equipe de entregar o sistema no prazo contra a performance necessária para lidar com o processamento de camadas vetoriais do mapa.
\end{regra}

\subsection{Opção A: Desempenho e Segurança (Rust + Dioxus)}
O uso da linguagem Rust com o framework Dioxus proporciona uma aplicação WebAssembly (Wasm) ou nativa extremamente otimizada, reduzindo o uso de memória no navegador cliente e aumentando consideravelmente a fluidez das transições de renderização vetorial.

\subsection{Opção B: Agilidade e Ecossistema (TanStack + React + Tailwind CSS)}
Devido à restrição de tempo, a stack baseada em React oferece a vantagem da vasta disponibilidade de componentes e bibliotecas prontas (ex: React Simple Maps, Leaflet). O gerenciamento de estados assíncronos seria tratado eficientemente com TanStack Query, enquanto a estilização ganharia agilidade através de utilitários do Tailwind CSS.

\section{Modelagem Preliminar de Dados}
A comunicação primária de persistência se dará via PostgreSQL, gerenciado pela aplicação Java.

\begin{entidade}{usuarios\_sistema}
\campo{id\_usuario}{uuid}{PK}
\campo{nome\_completo}{varchar(150)}{}
\campo{email}{varchar(100)}{UNIQUE}
\campo{senha\_hash}{varchar(255)}{}
\campo{nivel\_acesso}{varchar(20)}{}
\campo{data\_criacao}{timestamp}{}
\end{entidade}

\begin{entidade}{cache\_sinan}
\campo{id\_registro}{bigint}{PK}
\campo{codigo\_municipio}{varchar(10)}{FK (ibge)}
\campo{agravo}{varchar(100)}{}
\campo{competencia}{date}{}
\campo{total\_casos}{integer}{}
\end{entidade}

\vspace{1cm}
\begin{center}
\textit{Documento gerado como planejamento técnico preliminar para alinhamento da equipe de desenvolvimento.}
\end{center}

\end{document}

